\section{Introduction}

Accurate epidemiological characteristics are essential for evidence-based public health decision-making during infectious disease outbreaks.\cite{ref1,ref2,ref3,ref4} Parameters such as the basic reproduction number ($R_0$), serial interval, and case fatality rate serve as quantitative indicators of transmission potential, transmission speed, and disease severity, respectively. These parameters critically inform the development of public health strategies, the design of intervention measures, and the allocation of healthcare resources. However, obtaining reliable estimates of these parameters requires rigorous evidence synthesis.

Although parameter estimates from single high-impact studies or preliminary reports are commonly used, they often introduce substantial uncertainty and propagate systematic bias throughout model simulations, ultimately undermining the reliability and interpretability of projections. Systematic reviews and meta-analyses offer a superior alternative as the most authoritative source of such parameters. Through standardized integration of evidence across studies, they effectively reduce biases related to publication, geography, and population heterogeneity.\cite{ref5,ref6,ref7} Only by adopting parameters derived from such comprehensive evidence synthesis can transmission models move beyond theoretical constructs to deliver credible, policy-relevant predictions.

Nevertheless, despite generating high-quality evidence, the systematic review entails substantial resource requirements and considerable operational complexity. An empirical investigation of PROSPERO-registered projects has demonstrated that the average duration from protocol initiation to full publication of a systematic review extends to 67.3 weeks, about 1.29 years.\cite{ref8} Moreover, the literature screening process is particularly labor-intensive: initial database searches frequently retrieve tens of thousands of records, ranging from 27 to over 90,000, of which only about 2.94\% typically meet the prespecified eligibility criteria.\cite{ref8} Consequently, researchers must devote considerable effort to screening vast numbers of citations to identify a small set of eligible studies.

In response to this labor-intensive bottleneck, artificial intelligence has emerged as a promising solution. Recent advances in artificial intelligence (AI) have enabled AI-assisted tools to alleviate the burden of manual screening, with active learning algorithms demonstrating particular promise in prioritizing relevant citations while maintaining high sensitivity.\cite{ref9,ref10,ref11} Nevertheless, existing unmodified large language models remain incapable of independently executing the critical task of identifying and selecting relevant literature for systematic reviews. This limitation is starkly illustrated by the performance of three fully automated GPT-4-based approaches: Consensus GPT, Scholar GPT, and the native GPT web-browsing mode correctly identified merely 32.7\% (32/98), 19.4\% (19/98), and 6.1\% (6/98) of eligible articles, respectively.\cite{ref12} These poor results highlight the need for domain-specific adaptation. While specialized models such as MedCPT in biomedical information retrieval\cite{ref13} and TrialMind in clinical trial screening\cite{ref14} have shown substantially better performance, the field of AI-assisted systematic review still lacks a dedicated, high-performance retrieval framework.

Here, we present \metagent{}, a dimension-aware large language model framework developed for high-recall literature screening in systematic reviews of infectious disease modeling parameters. The framework evaluates each candidate article across five clinically and methodologically relevant dimensions: disease relevance, population relevance, location relevance, original empirical evidence, and parameter relevance; and employs a stage-aware screening strategy that adapts the input modality according to the availability of title and abstract information. By integrating these design elements, \metagent{} aims to preserve high sensitivity while substantially reducing the manual screening burden. This study demonstrates that dimension-aware LLMs can serve as effective assistive tools within human-in-the-loop systematic review workflows, providing a practical pathway to accelerate high-quality evidence synthesis for infectious disease modeling and public health decision-making.
